%% Language and font encodings
\usepackage[english]{babel}
% \usepackage[utf8x]{inputenc}
\usepackage[T1]{fontenc}


%% Sets page size and margins
% \usepackage[letterpaper,top=3cm,bottom=2cm,left=3cm,right=3cm,marginparwidth=1.75cm]{geometry}
\usepackage[margin=1.3in]{geometry}


%% Spacing
\usepackage{setspace}
% \doublespacing
 \onehalfspacing % ASI ESTABA 

%% Useful packages
\usepackage{amsmath}
\usepackage{amssymb}
\let\Bbbk\relax
\usepackage{newtxtext,newtxmath}
\usepackage{adjustbox}
\usepackage{graphicx}
\usepackage{rotating}

% \usepackage[colorinlistoftodos]{todonotes} activate for side notes

\usepackage[colorlinks=true, allcolors=blue]{hyperref}
\let\openbox\relax
\usepackage{amsthm}
\usepackage{lscape}
\usepackage{thmtools}
\usepackage{thm-restate}

\usepackage{color,soul}
\usepackage{xcolor}

\usepackage{tikz}

%% for tables
\usepackage{booktabs}
\usepackage{booktabs, caption}
\usepackage{tabularx}
\usepackage{array}
\newcolumntype{P}[1]{>{\centering\arraybackslash}p{#1}}
\newcolumntype{M}[1]{>{\centering\arraybackslash}m{#1}}
\usepackage[flushleft]{threeparttable}
\usepackage{multirow}
\usepackage{siunitx}
\usepackage{changepage}
%\usepackage{pdflscape}
% for placing figures
\usepackage{float}
\usepackage{placeins}
\usepackage{rotating} % For rotating tables
%\usepackage[nolists,tablesfirst]{endfloat}
%\usepackage{caption,rotating}
% citations
\usepackage{natbib}
% \bibliographystyle{abbrv}
%\bibliographystyle{abbrvnat}
\setcitestyle{authoryear,open={(},close={)}}
% \setlength{\bibsep}{0.0pt}

% \usepackage{bibspacing}
% \setlength{\bibitemsep}{.2\baselineskip plus .05\baselineskip minus .05\baselineskip}

% appendices
\usepackage[toc,page]{appendix}

% Custom definitions
\def\bmath#1{\mbox{\boldmath$#1$}}
\def\sym#1{\ifmmode^{#1}\else\(^{#1}\)\fi}
\declaretheorem[name=Theorem,numberwithin=section]{theorem}
\declaretheorem[name=Proposition,numberwithin=section]{proposition}
\declaretheorem[name=Lemma,numberwithin=section]{lemma}
\declaretheorem[name=Corollary,numberwithin=section]{corollary}
\declaretheorem[name=Definition,numberwithin=section]{definition}
\declaretheorem[name=Example,numberwithin=section]{example}


\newtheorem{assumption}{Assumption}
\newtheorem{property}{Property}

%for highlighting
\usepackage{soul}

%for including pdfs
\usepackage{pdfpages}

% \usepackage{titling}

% multiple footnotes
% \usepackage[multiple]{footmisc}
 
%  for cases in equations
 \usepackage{mathtools}
 
 % For comments

\newcounter{marginparcounter}
\newcommand{\paulapar}[1]{\stepcounter{marginparcounter}$^{(\roman{marginparcounter})}$\marginpar{\renewcommand{\baselinestretch}{0.8}\scriptsize$^{(\roman{marginparcounter})}$ PA: #1}}
\newcommand{\martinpar}[1]{\stepcounter{marginparcounter}$^{(\roman{marginparcounter})}$\marginpar{\renewcommand{\baselinestretch}{0.8}\scriptsize$^{(\roman{marginparcounter})}$ MS: #1}}
\newcommand{\kristianpar}[1]{\stepcounter{marginparcounter}$^{(\roman{marginparcounter})}$\marginpar{\renewcommand{\baselinestretch}{0.8}\scriptsize$^{(\roman{marginparcounter})}$ KLV: #1}}
 

\usepackage{graphicx} % for \includegraphics
\usepackage{subcaption} % for subfigure
\usepackage{rotating} % for sidewaystable and sidewaysfigure

%%%AEJ-Micro wants roman numerals for sections and capital letters for subsections
%\renewcommand{\thesection}{\Roman{section}}
\renewcommand\thesection{\Roman{section}}
\renewcommand\thesubsection{\Alph{subsection}}